\subsection{Materials}
During this experiment I used:
\begin{itemize}
    \item Holographic gratings $500$ lines/mm;
    \item holographic grating with unknown grating groove space;
    \item mobile phone;
    \item electric torch;
    \item bulb;
    \item measuring tape;
    \item support for mobile phone;
    \item support for electric torch.
\end{itemize}
The tape measure with a maximum range of $ 5 \, m $ and a sensitivity of $ 0.01 \, m $ was used to measure the distance between the light source and the mobile phone (observer), then to measure $ z $. In this experience I decided to use a flashlight (cold light) and the two bulb from my room lamp (warm light) in order to visualize possible differences in their chromatic components. To hold the flashlight and cell phone I built simple supports using Lego bricks in order to reduce the error associated with the $ z $ measurement and fix the source and the "observer" in a single position.
%Aggiungere immagine dei supporti
The experimental apparatus must be prepared in a sufficiently large space. If the distance $ z $ is too small there is the risk that the first order of diffraction does not appear on the second holographic grating, making the measurement of its step not possible. 

The source must be placed at a distance $ z $ from the sensor with the holographic grating in front of it. It's important to place the source and the sensor at the same height.
%Aggiungere schema apparato

\subsection{Procedures}
After having prepared the experimental setup, the first step consists in measuring the distance $ z $ that separates the sensor (with the holographic grating in front) and the source. After, through the first holographic grating with a known $d$, the real experiment can start. As mentioned in the previous section I decided to analyze two different sources but both are shared by the data taking method. Once everything was positioned, through the sensor I took images with and without the holografic grating without moving the sensor or the source. I took the images without the reticle with the light on in order to better visualize the support of the source that I take as a measurement reference. I took the images with the holographic reticle with all the lights in the room off and setting a low exposure time from the mobile phone in order to capture the chromatic components at their best.

After obtaining all the photos with both holographic gratings, I moved the photos to the PC and downloaded Imagej software to analyze and measure the distances between the source and its chromatic components. Using the source support measured previously and clearly visible in the photos taken with the light on as a reference measure, I was able to measure the distance between the source and the first order of diffraction. For the first holographic grating, measuring the distance $\Delta x$ is used to obtain the wavelength $\lambda$ using the formula (\ref{deltax}). Subsequently with the same procedure using the Imagej software I measure again the $\Delta x $ from the image of the second holographic grating with $d$ unknown and through $\lambda$ previously measured and $\Delta x$ with the formula (\ref{d}) I obtain $d'$.

